\uuid{7vYj}
\exo7id{831}
\titre{exo7 831}
\auteur{ridde}
\organisation{exo7}
\datecreate{1999-11-01}
\isIndication{false}
\isCorrection{false}
\chapitre{Calcul d'intégrales}
\sousChapitre{Intégrale abélienne}
\cpge{mpsi-mp2i.primitives-equations-differentielles.a, pcsi.techniques-calcul-primitives.a, ptsi.techniques-calcul-primitives.a, mpsi-mp2i.fonctions-valeurs-reelles-complexes.b, pcsi.techniques-calcul-fonctions.b, ptsi.techniques-calcul-fonctions.b, mpsi-mp2i.polynomes-fractions-rationnelles.h, bcpst1.derivees-primitives-integrales.d, bcpst1.derivees-primitives-integrales.a}
\module{Analyse}
\niveau{L1}
\difficulte{}

\contenu{
\question{
D\'eterminer les intervalles d'\'etude et calculer les primitives des fonctions :
$$ \frac{8x-3}{\sqrt{12x-4x^2-5}}$$
$$ \sqrt{x^2-1}$$
$$ \frac{x \sqrt x}{x^2-5x + 4}$$
}
}
